\section{A residue-aware Cauchy separation}
\label{sec:pivot21}

This section changes the separated basis used for crossing entries in one
contraction call. It does not assume an unproved faster summation routine.
The new basis has the same source and target groups as the previous
factorial-Cauchy basis, but its guaranteed valuation grows with slope
$64/21$ instead of $3$.

\subsection{The actual row domain}

Let $\mathcal O$ be the unramified coefficient ring used in the source
construction, with valuation $\nu_2(2)=1$. In particular, its reduction
modulo two can be a proper extension of $\mathbb F_2$. Set
\[
 L=2^a,\quad i=\nu_2(t)<a,\quad
 \alpha=2^{a+1},\quad T=t+\alpha\zeta,\quad \zeta\in\mathcal O,
 \qquad J=2^{a-i-1}.
\]
A column tile has positions
$y=y_0+2^{i+1}j$, where $0\le y_0<2^{i+1}$ and $0\le j<J$.
For a crossing row $x$, the bits in position $i$ of $x$ and $y_0$
are equal. Define
\[
 d_0=t-y_0+x,\qquad d=d_0+\alpha\zeta,\qquad
 z_x=\frac{2^{i+1}}{d},\qquad w_x=\frac{z_x}{2},
 \qquad R=a+1-i.
\]
The source's crossing condition gives $d_0\ne0$ and
$\nu_2(d_0)\le i$. With
\[
 \mathscr C(T)=\prod_{q=-L}^{L-1}(T-q),\qquad
 c_x=\frac{\mathscr C(T)}{\mathscr C'(0)d},
\]
the kernel is
\begin{equation}
 W_{xy}(T)=\frac{c_x}{1-z_xj},
 \qquad z_x\in2\mathcal O,\quad \nu_2(c_x)\ge1.
 \label{eq:pivot21-kernel}
\end{equation}
The last lower bound may be sharpened to $R$, but an offset of one
already suffices for the comparison below.

\begin{lemma}[A missing residue class]
\label{lem:pivot21-domain}
Put $w_0=2^i/d_0\in\mathbb Z_{(2)}$. Then
\[
 w_x-w_0\in2^R\mathcal O.
\]
If $i\ge1$, define $c\in\{2,6\}$ by
\[
 c=\begin{cases}
       6,&\text{bit }i-1\text{ of }y_0\text{ is zero},\\
       2,&\text{bit }i-1\text{ of }y_0\text{ is one}.
    \end{cases}
\]
For $R\ge3$, the residue of $w_x$ modulo eight lies in the prime
subring and is different from $8-c$. If $i=0$, $w_x$ is odd and
either choice of $c$ is valid.
\end{lemma}

\begin{proof}
Write $v=\nu_2(d_0)\le i$. Adding $2^{a+1}\zeta$ preserves that
valuation. Therefore
\[
 w_x-w_0=
 -\frac{2^{i+a+1}\zeta}{d_0(d_0+2^{a+1}\zeta)},\qquad
 \nu_2(w_x-w_0)\ge i+a+1-2v\ge R.
\]
This proves the prime-subring assertion when $R\ge3$.

Only the case $\nu_2(w_0)=1$ could give residue two or six. In this
case $v=i-1$. Write the lower $i$ bits of $x$ and $y_0$ as $x_<$
and $y_<$. Since their bits in position $i$ agree, the other part of
$x-y_0$ is divisible by $2^{i+1}$. The difference $x_<-y_<$ lies
strictly between $-2^i$ and $2^i$ and, by $v=i-1$, is an odd
multiple of $2^{i-1}$. It must be $+2^{i-1}$ when bit $i-1$ of
$y_0$ is zero, and $-2^{i-1}$ otherwise. Since
$t/2^{i-1}\equiv2\pmod4$, the odd number
$d_0/2^{i-1}$ is respectively three or one modulo four. Its inverse
has the same residue, so $w_0$ is respectively six or two modulo
eight. In all other valuation cases neither residue can occur.
For $i=0$, $d_0$ is odd, which proves the last statement.
\end{proof}

The exceptional value $R=2$ gives $J=1$. It has only the mode $b=0$
and needs no residue exclusion.

\subsection{An explicit sequence of auxiliary row points}

Repeat the word of length 21
\[
 A\ B\ A\ B\ C\ A\ B\ A\ B\ C\ A\ B\ A\ A\ C\ A\ B\ A\ A\ C\ A.
\]
At the $k$th occurrence of a letter, starting at $k=0$, set
\begin{equation}
 v_h=\begin{cases}
       1+2k,&\text{letter }A,\\
       4k,&\text{letter }B,\\
       c+8k,&\text{letter }C,
     \end{cases}
 \qquad \rho_h=2v_h.
 \label{eq:pivot21-points}
\end{equation}
The index $h$ counts positions in the repeated word. Let $A_b,B_b,C_b$
be the numbers of the three letters among its first $b$ positions.
Thus a complete block has counts $(11,6,4)$. Define
\begin{equation}
 \eta_b=\min\left\{
 A_b+\nu_2(A_b!),\quad
 2B_b+\nu_2(B_b!)+C_b,\quad
 3C_b+\nu_2(C_b!)+B_b
 \right\}.
 \label{eq:pivot21-eta}
\end{equation}

\begin{lemma}[Factorials near the prime subring]
\label{lem:pivot21-near-prime}
If $q\in\mathcal O$, $q_0\in\mathbb Z_{(2)}$, and
$q-q_0\in2^Q\mathcal O$, then for every integer $0\le m<2^Q$,
\[
 \nu_2\!\left(\prod_{k=0}^{m-1}(q-k)\right)\ge\nu_2(m!).
\]
\end{lemma}

\begin{proof}
For $1\le r\le\lfloor\log_2m\rfloor$, the rational $q_0$ has a
well-defined integer residue modulo $2^r$. Among $m$ consecutive
integers, at least $\lfloor m/2^r\rfloor$ have that residue.
Since $r\le Q$, the corresponding factors $q-k$ are divisible by
$2^r$. Summing over $r$ proves the result by Legendre's formula.
For $m=0,1$ the assertion is immediate.
\end{proof}

This lemma is necessary: an arbitrary extension element of $\mathcal O$
need not have integral binomial coefficients. The actual rows have the
required proximity to the prime subring.

\begin{lemma}[Valuation of the pivot product]
\label{lem:pivot21-product}
For every actual crossing row and every $0\le b<J$,
\[
 \nu_2\!\left(\prod_{h=0}^{b-1}(w_x-v_h)\right)\ge\eta_b.
\]
Furthermore, $\eta_b$ is nondecreasing and
\[
 \eta_b=\frac{22}{21}b+O(\log(b+2)).
\]
\end{lemma}

\begin{proof}
Suppose first that $R\ge3$. The permitted residues split into three
cases. If $w_x$ is odd, the $A$ factors give
\[
 2^{A_b}\prod_{k<A_b}\left(\frac{w_x-1}{2}-k\right),
\]
while every $B$ or $C$ factor is a unit. Their valuation is at least
$A_b+\nu_2(A_b!)$.

If $w_x\equiv0\pmod4$, the $B$ factors similarly contribute at
least $2B_b+\nu_2(B_b!)$, each $C$ factor contributes exactly one
power of two, and the $A$ factors are units. The resulting bound is
the second term in \eqref{eq:pivot21-eta}.

The remaining permitted case is $w_x\equiv c\pmod8$. Now the $C$
factors contribute at least $3C_b+\nu_2(C_b!)$, each $B$ factor
contributes one power of two, and the $A$ factors are units. This
gives the third term.

To justify the factorial steps over $\mathcal O$, apply
Lemma~\ref{lem:pivot21-near-prime} after dividing by two, four, or
eight. Lemma~\ref{lem:pivot21-domain} supplies precisions $R-1$,
$R-2$, and $R-3$ respectively. The first two factor counts satisfy
\[
 A_b\le b<J=2^{R-2}<2^{R-1},\qquad
 B_b\le b<2^{R-2}.
\]
If $b=21q+r$, $0\le r\le20$, the positions of the letter $C$ give
\[
 C_b=4q+\lfloor r/5\rfloor\le\lfloor b/5\rfloor
 \le\lfloor b/4\rfloor.
\]
Thus, whenever $C_b>0$, $C_b<2^{R-3}$. If $R\le3$, no $C$ factor
occurs. The case $R=2$ has $b=0$, so its product bound is immediate.

Each expression in the minimum \eqref{eq:pivot21-eta} is
nondecreasing as the three counts increase. For $b=21q$,
Legendre's formula gives the exact identity
\[
 \eta_{21q}=22q-
 \max\{s_2(11q),s_2(6q),s_2(4q)\},
\]
where $s_2$ is the number of ones in the binary expansion.
For general $b$, the three counts differ by $O(1)$ from
$(11b/21,6b/21,4b/21)$. The bound
$s_2(m)=O(\log(m+2))$ proves the claimed asymptotic formula.
\end{proof}

\subsection{A finite separation with no nonunit division}

For arbitrary even auxiliary points $\rho_h$, set
\[
 U_b(z)=\prod_{h<b}\frac{z-\rho_h}{1-hz},\qquad
 V_b(j)=\prod_{h<b}\frac{j-h}{1-\rho_hj}.
\]
All denominators are units for $z\in2\mathcal O$ and integer $j$.

\begin{lemma}[Finite Cauchy elimination]
\label{lem:pivot21-finite}
For every nonnegative integer $q$,
\begin{equation}
 \frac1{1-zj}
 =\sum_{b<q}
 \frac{(1-b\rho_b)U_b(z)V_b(j)}{(1-bz)(1-\rho_bj)}
 +\frac{U_q(z)V_q(j)}{1-zj}.
 \label{eq:pivot21-finite}
\end{equation}
If $0\le j<J$ and $q=J$, the remainder is zero.
\end{lemma}

\begin{proof}
The elementary identity
\[
 (1-bz)(1-\rho_bj)-(z-\rho_b)(j-b)
   =(1-b\rho_b)(1-zj)
\]
shows that the $b$th summand is the difference of the consecutive
remainders $U_b(z)V_b(j)/(1-zj)$ and
$U_{b+1}(z)V_{b+1}(j)/(1-zj)$. Telescope. For $j<J$, the product
$V_J(j)$ contains the zero factor $j-j$.
\end{proof}

For the specific points \eqref{eq:pivot21-points}, define
\begin{align}
 \ell_b(j)&=
 \frac{(1-b\rho_b)\binom jb}
      {\prod_{h=0}^{b}(1-\rho_hj)},
 \label{eq:pivot21-source}\\
 r_b(x)&=
 c_x\,b!\,
 \frac{\prod_{h=0}^{b-1}(z_x-\rho_h)}
      {\prod_{h=0}^{b}(1-hz_x)}.
 \label{eq:pivot21-target}
\end{align}
Here $\binom jb=0$ for $b>j$.

\begin{theorem}[Residue-aware separated kernel]
\label{thm:pivot21-kernel}
For every actual crossing entry,
\[
 W_{xy}(T)=\sum_{b=0}^{J-1}r_b(x)\ell_b(j).
\]
Every source feature $\ell_b(j)$ belongs to $\mathbb Z_{(2)}$.
The target coefficient satisfies
\begin{equation}
 \nu_2(r_b(x))\ge\omega_b,
 \qquad
 \omega_b=1+b+\nu_2(b!)+\eta_b
          =\frac{64}{21}b+O(\log(b+2)).
 \label{eq:pivot21-weight}
\end{equation}
The weights $\omega_b$ are strictly increasing. The stronger offset
$R$ may replace one in this formula.
\end{theorem}

\begin{proof}
Use Lemma~\ref{lem:pivot21-finite} with $q=J$, and replace
$\prod_{h<b}(j-h)$ by the integer product $b!\binom jb$.
The integer binomial and the odd denominators make
\eqref{eq:pivot21-source} integral. In
\eqref{eq:pivot21-target}, factoring two from each of the $b$
numerator factors gives $2^b\prod_{h<b}(w_x-v_h)$. Combine
Lemma~\ref{lem:pivot21-product} with the explicit $b!$ and
$\nu_2(c_x)\ge1$. Legendre's formula gives the slope
$1+1+22/21=64/21$. Finally, $\eta_b$ is nondecreasing and the term
$b$ increases strictly, so the weights do also.
\end{proof}

The formula multiplies the factorial into the target. It never divides
by a factorial in the coefficient ring. Implementation uses integer
binomial coefficients, ring multiplication, and unit inversion only.
The points $\rho_h$ are auxiliary interpolation points; they need not
be actual rows in a packet.

\subsection{Consequences for the contraction accounting}

For each existing column tile, replace its source moment of degree $b$
by
\[
 L_b(X)=\sum_{\substack{\text{real columns }y\text{ in the tile}\\
                        y=y_0+2^{i+1}j}}
               \ell_b(j)X_y.
\]
There is still one mode of each weight per old tile. It is a linear
form in bits of the same source group and targets the same other group.
After expanding the row factors, a mode-exponent tuple $\beta$ whose
total certified weight is at least the working precision $p$ has
coefficient zero modulo $2^p$. Thus the retained condition is
\[
 \sum_{\text{tile families }f}\sum_{b=0}^{J-1}
             \omega_b\beta_{f,b}<p.
\]
The source's child factorization, wire count, degree bounds, and
creation-precision convention are preserved. Products defining the
features contain at most $J\le L$ factors of polynomial description
length, so the entry-oracle bounds remain valid. This argument also
permits the strongest compatible previous packet-entry routine.

To quantify the gain, let $N_\gamma(c,p)$ denote the abstract tuple
count with $c$ families and weights
$\omega_b=\gamma b+O(\log(b+2))$. Its standard generating-product
bound is
\[
 N_\gamma(c,p)\le
 \exp\!\left(\theta p+
       c\sum_{b\ge0}-\ln(1-e^{-\theta\omega_b})\right).
\]
As $\theta\downarrow0$, the sum has leading term
$\pi^2/(6\gamma\theta)$. Optimizing $\theta$ gives the leading
logarithmic bound
\[
 \ln N_\gamma(c,p)\le
 \left(\pi\sqrt{\frac{2}{3\gamma}}+o(1)\right)\sqrt{cp}.
\]
The actual finite cutoff only reduces this bound. In the active-call
regime, $p/c\to\infty$, so the displayed asymptotic applies uniformly
with the previously verified schedule.

Consequently, replacing the previous factorial-Cauchy weights of slope
$3$ by \eqref{eq:pivot21-weight}, while retaining a matched precision
and recursion schedule, multiplies the leading enumeration coefficient by
\begin{equation}
 \sqrt{\frac{3}{64/21}}=\sqrt{\frac{63}{64}}
   =0.9921567416\ldots.
 \label{eq:pivot21-ratio}
\end{equation}
Relative to the geometric weights of slope one, the factor is
$\sqrt{21/64}$. These are comparisons of explicit upper-bound
certificates. They are not wall-clock speedup measurements or lower
bounds on the old algorithm, and they do not change the logarithmic
order of the global SAT bound.

The new weights need not dominate the old weights at every finite
index. For example, the new weight at $b=4$ is $10$, while the old
factorial-Cauchy weight is $11$. The improvement is in the proved
leading asymptotic coefficient. If desired, an implementation may
count the old and new abstract tuple sets by ordinary generating-function
dynamic programming and choose the smaller set at each call. The numbers
of families, modes and precision units are polynomial in $s$, and the
counts have polynomial bit length, so this selection has polynomial
bit cost and space. It preserves both exact identities and the new
asymptotic upper bound. No uniform finite-input speedup is asserted.
